% EMNLP/ARR paper source with switchable modes.
% Set \anonymoussubmissiontrue for ARR anonymous submission.
% Set \preprintversiontrue for arXiv/preprint builds.
\documentclass[11pt,a4paper]{article}
\newif\ifanonymoussubmission
\newif\ifpreprintversion
\anonymoussubmissiontrue
\preprintversionfalse
% Remap fonts unavailable in this TeX installation to CM equivalents
\renewcommand\sfdefault{cmss}
\renewcommand\ttdefault{cmtt}
\ifanonymoussubmission
  \usepackage[review]{acl}
\else
  \ifpreprintversion
    \usepackage[preprint]{acl}
  \else
    \usepackage{acl}
  \fi
\fi

% Standard package includes
\usepackage{times}
\renewcommand\ttdefault{cmtt}  % times sets ttdefault to pcr (Courier), not available here
\usepackage{latexsym}
\usepackage[T1]{fontenc}
\usepackage[utf8]{inputenc}
\usepackage{microtype}
\usepackage{amsmath}
\usepackage{amssymb}
\usepackage{graphicx}
\usepackage{listings}
\lstset{
  basicstyle=\ttfamily\footnotesize,
  breaklines=true,
  breakatwhitespace=false,
  columns=flexible,
  keepspaces=true,
  frame=single,
  xleftmargin=4pt,
  xrightmargin=4pt,
}
\usepackage{booktabs}
\usepackage{multirow}
\usepackage{multicol}
\usepackage{array}

\ifanonymoussubmission
\setlength\linenumbersep{4pt}
\renewcommand\linenumberfont{\normalfont\tiny\color{gray}}
\AtBeginDocument{\leftlinenumbers*}
\fi

\title{Accuracy, Stability, and Repeated-Run Reliability of Large Language Models on Deterministic Programming Tasks}

\ifanonymoussubmission
\author{Anonymous ACL Submission}
\else
\author{
  Yongxi Zhou \\
  Northeastern University \\
  Massachusetts, USA \\
  \texttt{zhou.yongx@northeastern.edu}
}
\fi

\begin{document}
\maketitle

\begin{abstract}
Run-level pass rate overstates retry-free coverage by up to 17.8 percentage points---and the gap is largest precisely for mid-performing systems. We investigate this accuracy--stability relationship in large language model (LLM) evaluation for deterministic text-conditioned generation, using programming tasks as a concrete testbed. Standard code-generation benchmarks emphasize single-run accuracy or eventual success under repeated sampling, but many deployment settings also require \emph{stability}: consistent outcomes across repeated invocations under the same task description. We present a repeated-run evaluation protocol with metrics for run-level accuracy, retry-free coverage, and per-problem variability. On a recency-based benchmark of 100 LeetCode-style problems, we evaluate 16 models from five provider families under two prompt templates with five repeated runs per problem, yielding 16{,}000 evaluation instances. Although run-level pass rate and perfect stability rate are strongly correlated ($r=0.985$), pass rate consistently exceeds retry-free coverage---a gap that reaches 17.8 percentage points and reverses model rankings even among closely matched systems. Prompt effects are model-dependent rather than uniformly beneficial. These results suggest that repeated-run stability analysis is a necessary complement to conventional accuracy reporting for deterministic text-conditioned generation tasks.
\end{abstract}

\section{Introduction}
Large language models are increasingly used for text-conditioned tasks with deterministic specifications: given a fixed textual input and a fixed verifier, correctness is binary. Program synthesis and code assistance are especially important instances of this broader setting, because they pair natural-language problem statements with executable test suites. Conventional benchmarks typically report single-run performance, or summarize repeated sampling through eventual-success metrics such as pass@$k$ \cite{chen2021codex,austin2021program}. Recent work has shown that repeated sampling can be scaled as a form of inference-time compute \cite{brown2024large}, yet the \emph{stability} of outcomes across such repeated invocations---whether a model consistently succeeds or fails on the same task---has received less systematic attention. In practice, repeated calls under the same task description can yield different outputs and different pass/fail outcomes, even when the downstream workflow expects reproducible behavior \cite{atil2024llm,alvarado2025repetitions}.

This paper studies that repeated-invocation variability directly. Rather than treating repeated sampling only as an opportunity to recover one success among many attempts, we investigate the relationship between correctness and stability under fixed conditions. Our focus is therefore complementary to pass@$k$ \cite{chen2021codex}: whereas pass@$k$ asks whether a model can produce \emph{at least one} correct solution in $k$ attempts---a capability measure---we ask how reliably it succeeds on \emph{every} invocation, a deployment reliability measure. These two questions target different downstream concerns and can yield divergent rankings, particularly for mid-accuracy systems.

The paper is organized around three empirical questions:
\begin{enumerate}
  \item How large is the gap between run-level success and retry-free problem coverage under repeated evaluation?
  \item Which aggregate summaries best expose problem-level variability that single-run reporting hides?
  \item As an exploratory analysis: do prompt choices or reasoning-specialized architectures systematically shift repeated-run stability?
\end{enumerate}

We focus on deterministic tasks because they represent a broader and practically important class of NLP system behavior: the model is conditioned on text, but the output is consumed by a verifier that expects reproducible correctness. Code-generation systems are a particularly stringent instance of this pattern, alongside tasks such as text-to-SQL \cite{gao2023text}, structured semantic parsing, tool-call argument generation, and formal text transformations. In such settings, the accuracy--stability gap has direct practical consequences: a model with 60\% run-level pass rate but only 47\% perfect stability solves many individual attempts, yet still leaves a substantial fraction of tasks in a retry-required regime. Our aim is therefore not merely to rank models by mean performance, but to characterize how much of their observed accuracy translates into reliable repeated-use behavior.

\section{Related Work}
Code-generation evaluation is commonly framed around functional correctness and pass@$k$ under automated testing \cite{chen2021codex,austin2021program}. Foundational benchmarks emphasize whether a model eventually produces at least one correct program under repeated sampling \cite{chen2021codex,austin2021program,hendrycks2021apps}, while more recent efforts address contamination and temporal validity \cite{jain2024livecodebench}, with parallel work extending evaluation to real-world repository-level issues \cite{jimenez2024swebench}. Prompting and inference-time interventions use additional computation to improve performance rather than to measure it: zero-shot chain-of-thought \cite{kojima2022zeroshot} restructures the prompt, while self-consistency and adaptive-consistency \cite{wang2022self,aggarwal2023lets} aggregate multiple samples. In parallel, work on LLM output non-determinism shows that performance varies substantially across repeated runs under identical conditions \cite{atil2024llm,alvarado2025repetitions}, and prompt sensitivity research demonstrates that small phrasing changes can shift accuracy \cite{zhuo2024prosa}.

Perfect Stability Rate (PSR) is mathematically equivalent to pass@$k$ at $k=R$ (the number of repeated runs) with an all-successes threshold; the distinction is interpretive. Pass@$k$ estimates $P(\text{at least one success in }k\text{ draws})$---a \emph{capability} measure---while PSR estimates $P(\text{all }k\text{ draws succeed})$---a \emph{deployment reliability} measure. These quantities diverge most sharply for mid-accuracy systems. The contribution lies in treating PSR and Average Variance (AV) as first-class evaluation targets for deterministic code tasks---both defined formally in Section~\ref{sec:metrics}---and using them to characterize the accuracy--stability gap directly.

Recent work emphasizes that code evaluation is sensitive to benchmark construction and reporting choices \cite{liu2024your,du2024evaluating}. Our framing shifts the emphasis from benchmark validity to repeated-invocation reliability under a fixed judge. Work on confidence calibration \cite{guo2017calibration} and performance--efficiency trade-offs \cite{zhang2026performance} reflects a broader recognition that average accuracy alone is an incomplete measure of model reliability. This perspective is deliberately narrower than end-to-end agent evaluation: it isolates repeated-run behavior for single-turn coding, while broader orchestration, tool use, and runtime dependencies require separate analysis.

\section{Experimental Setup}

\subsection{Problem Dataset}
We evaluate on a curated dataset of deterministic programming problems, denoted $\mathcal{D}=\{(x_i, T_i)\}_{i=1}^{N}$, where $x_i$ is the natural-language specification (including I/O format and constraints) and $T_i$ is a deterministic test harness. Each problem is graded by executing the model-produced program against $T_i$, producing a binary correctness indicator. These tasks are instantiated as LeetCode problems with platform-provided acceptance criteria, and all generations are evaluated as Python solutions against Python-based submission templates.

\paragraph{Problem selection.} To reduce the risk of benchmark contamination from model training data, we select the $N=100$ most recently published problems available on the platform at collection time, sampling the newest problems first. Problems requiring paid-only access are excluded. No further manual curation is applied; the final benchmark is entirely determined by recency rank. This construction does not guarantee that evaluated models have never encountered related problem statements during training, but it substantially reduces the probability of direct memorization compared to historically popular problems. We note, however, that recency alone does not guarantee a contamination-free evaluation: recently published contest problems are often discussed publicly online---including on forums, social media, and editorial sites---within days of release. Post-release discussion means that problem statements, constraints, and accepted solutions may appear in model training corpora even for recently added problems, and we do not perform contamination-detection analysis (e.g., via prompt perturbation, n-gram overlap, or membership inference \cite{deng2023investigating}). This limitation is discussed further in the Limitations section.

The dataset spans three difficulty tiers: 20 Easy, 50 Medium, and 30 Hard problems, covering a broad range of algorithmic topics including Array, Dynamic Programming, String, Graph, and Simulation. All problems satisfy the following properties:
\begin{itemize}
  \item \textbf{Deterministic grading:} each $T_i$ yields a deterministic accept/reject outcome.
  \item \textbf{Self-contained evaluation:} no network access or external dependencies are required at runtime.
  \item \textbf{Fixed snapshot:} all models are evaluated on the identical problem set to ensure comparability across runs.
  \item \textbf{No personal data:} problems consist of algorithmic specifications only and contain no personally identifying information.
\end{itemize}

In the evaluation framework, this benchmark is stored as a fixed local dataset snapshot rather than fetched online at run time. This design ensures that all model families are evaluated against the same problem set and ordering, and reduces the chance that later platform changes silently alter the benchmark during the experimental campaign.

\subsection{Models}

We evaluate 16 models across five provider families, covering a range of model scales, training objectives, and architectural choices, including both standard and reasoning-specialized variants. The evaluated models are listed in Table~\ref{tab:models}.

\begin{table}[t]
\centering
\small
\caption{Models evaluated in this study, grouped by provider family. $\dagger$ denotes reasoning-specialized models. Display names in this table are editorial labels derived from API identifier strings; they are not official product version names. Exact API identifiers and access dates are listed in Appendix~\ref{sec:model-snapshots}. Google models marked ``(preview)'' were accessed via preview-tier API endpoints.}
\label{tab:models}
\resizebox{\linewidth}{!}{
\begin{tabular}{lll}
\toprule
Family & Model & Type \\
\midrule
\multirow{3}{*}{OpenAI}
  & GPT-4.1         & flagship \\
  & GPT-4.1-mini    & lightweight \\
  & o4-mini         & reasoning$^\dagger$ \\
\midrule
\multirow{2}{*}{Anthropic}
  & Claude Sonnet 4.5 & flagship \\
  & Claude Haiku 4.5  & lightweight \\
\midrule
\multirow{5}{*}{Google}
  & Gemini 3.1 Pro (preview)        & flagship \\
  & Gemini 3 Flash (preview)        & efficient \\
  & Gemini 3.1 Flash-Lite (preview) & lightweight \\
  & Gemini 2.5 Flash                & efficient \\
  & Gemini 2.5 Flash+Think          & reasoning$^\dagger$ \\
\midrule
\multirow{4}{*}{Qwen}
  & QwQ-Plus   & reasoning$^\dagger$ \\
  & Qwen-Plus  & flagship \\
  & Qwen-Max   & large \\
  & Qwen-Turbo & lightweight \\
\midrule
\multirow{2}{*}{DeepSeek}
  & DeepSeek-R1   & reasoning$^\dagger$ \\
  & DeepSeek-V3.2 & flagship \\
\bottomrule
\end{tabular}
}
\end{table}

\subsection{Evaluation Framework}
For a model $m$ and a fixed evaluation configuration $c$ (prompt template, decoding parameters, toolchain), we run each problem $x_i$ multiple times. Each run produces a candidate program which is executed against $T_i$.

Let $Y_{i,r}^{(m,c)} \in \{0,1\}$ denote whether run $r \in \{1,\dots,R\}$ passes all tests for problem $i$. We treat compilation errors, runtime errors, timeouts, and wrong answers as failures ($0$).

The evaluation pipeline separates generation from execution. For large-scale experiments, model outputs are first generated through provider APIs, cached locally, normalized into a shared format, and only then submitted to the deterministic judge. This design makes experiments resumable across providers and reduces the risk that transient provider-side failures are conflated with downstream execution outcomes.

\subsection{Repeated-Run Methodology}
We adopt a repeated-run design to estimate both mean performance and variability. For each $(m,c)$ pair we:
\begin{enumerate}
  \item Sample $R$ independent generations per problem under configuration $c$.
  \item Evaluate each generation deterministically against $T_i$.
  \item Aggregate outcomes into run-level and problem-level metrics with uncertainty estimates.
\end{enumerate}

The primary controlled intervention is the prompt template. We compare two prompt variants while holding the dataset, model, decoding parameters, and execution environment fixed. All models are evaluated with temperature $T=0.3$, top-$p=0.9$, and $R=5$ repeated runs per problem per prompt configuration. The default maximum output budget is $4{,}096$ tokens; reasoning-specialized models (o4-mini, Gemini 2.5 Flash+Think) use larger budgets to accommodate extended chain-of-thought output (see Appendix~\ref{sec:model-snapshots} for per-model details). The o-series API (o4-mini) does not support explicit temperature or top-$p$ settings; decoding parameters are omitted for that model (implications discussed in the Limitations section). We use a non-zero temperature because the goal of the study is to characterize repeated-run variability under realistic stochastic decoding rather than to collapse outputs toward near-deterministic behavior at $T=0$. The two prompt templates---\textsc{Detailed} and \textsc{Minimal}---are held fixed across all models to enable a controlled comparison of prompt sensitivity.

\paragraph{Prompt templates.} Both prompts instruct the model to return plain Python source code only and to follow the provided LeetCode method signature exactly. The \textsc{Detailed} prompt uses a longer structured template emphasizing careful constraint reading, edge-case coverage, submission readiness, and accepted-style performance. The \textsc{Minimal} prompt conveys the same task more compactly, with lighter instruction scaffolding. Because both templates share the same output constraints and code template, the prompt comparison is a controlled low-contrast comparison that isolates instruction density rather than changing the task itself. Accordingly, the study does not test broad prompt-engineering strategies; it evaluates a narrow prompt variation under otherwise fixed conditions.

\paragraph{Normalization and execution.} Raw generations from different providers are normalized into a common local representation before evaluation. Code is extracted from model responses, basic syntax validity is checked, and each candidate is then submitted to the same deterministic judge. This normalization step is important because provider APIs differ in output packaging, token accounting, and batch interfaces even when the downstream task is identical.

\subsection{Metrics Definition}
\label{sec:metrics}

\paragraph{Run-Level Pass Rate (RLPR).}
$$
\mathrm{RLPR}(m,c)=\frac{1}{NR}\sum_{i=1}^{N}\sum_{r=1}^{R} Y_{i,r}^{(m,c)}.
$$
RLPR measures the probability that a \emph{random invocation} succeeds.

\paragraph{Perfect Stability Rate (PSR).}
$$
\mathrm{PSR}(m,c)=\frac{1}{N}\sum_{i=1}^{N} \mathbf{1}\!\left[\sum_{r=1}^{R} Y_{i,r}^{(m,c)}=R\right].
$$
PSR quantifies the fraction of tasks for which the system is reliably correct without retries.

\paragraph{Average Variance (AV).}
$\mathrm{AV}(m,c)=\frac{1}{N}\sum_{i=1}^{N} \frac{R}{R-1}\hat{p}_i(1-\hat{p}_i)$ with $\hat{p}_i=\frac{1}{R}\sum_r Y_{i,r}$ captures average instability across tasks.

\paragraph{Inference.}
We report 95\% Wilson score intervals for RLPR and PSR, and nonparametric bootstrap CIs for AV. Since both prompt templates are evaluated on the same problems, we compare PSR across prompt conditions using exact McNemar tests on the binary indicator of whether all $R=5$ runs pass. Reported statistics are descriptive summaries of the observed protocol rather than universal claims; prompt comparisons are specific to the two templates used here.

\section{Results}

\begin{table*}[t]
\centering
\small
\caption{Main results across all 16 models and two prompt configurations (100 problems, $R=5$ runs, 16{,}000 total evaluation instances). RLPR and PSR are in \%; AV is shown $\times 10^{-2}$ (i.e., multiply by 0.01 for the raw value). All entries are point estimates; 95\% CIs (Wilson score for RLPR/PSR; bootstrap for AV) available from authors on request. $\dagger$ reasoning-specialized model. $*$ output budget differs from the 4{,}096-token default: Gemini 2.5 Flash+Think used 24{,}576 tokens (8{,}192 thinking); o4-mini used 16{,}384 tokens. $\ddagger$ o4-mini decoding is not user-controllable; its PSR/AV may reflect reduced stochasticity rather than task reliability---see the Limitations section.}
\label{tab:overall-performance}
\setlength{\tabcolsep}{2.5pt}
\resizebox{\textwidth}{!}{
\begin{tabular}{llcccccc}
\toprule
\multirow{2}{*}{Family} & \multirow{2}{*}{Model}
  & \multicolumn{2}{c}{RLPR (\%)}
  & \multicolumn{2}{c}{PSR (\%)}
  & \multicolumn{2}{c}{AV ($\times 10^{-2}$)} \\
\cmidrule(lr){3-4}\cmidrule(lr){5-6}\cmidrule(lr){7-8}
& & D & M & D & M & D & M \\
\midrule
\multirow{3}{*}{OpenAI}
  & GPT-4.1               & 59.6 & 59.2 & 47.5 & 49.0 & 4.36 & 4.08 \\
  & GPT-4.1-mini          & 58.4 & 59.2 & 50.0 & 47.0 & 3.28 & 4.48 \\
  & o4-mini$^{\dagger\ddagger}$ & 86.2 & 85.6 & 75.0 & 69.0 & 3.36 & 5.20 \\
\midrule
\multirow{2}{*}{Anthropic}
  & Claude Sonnet 4.5 & 52.8 & 51.4 & 43.0 & 47.0 & 3.36 & 1.68 \\
  & Claude Haiku 4.5  & 41.2 & 40.8 & 35.0 & 34.0 & 2.40 & 2.00 \\
\midrule
\multirow{5}{*}{Google}
  & Gemini 3.1 Pro (preview)        & \textbf{86.2} & 85.8 & \textbf{75.0} & 74.0 & 3.36 & 3.92 \\
  & Gemini 3 Flash (preview)        & 80.8 & \textbf{83.2} & 67.0 & 70.0 & 5.28 & 4.40 \\
  & Gemini 3.1 Flash-Lite (preview) & 69.2 & 63.6 & 53.0 & 50.0 & 5.28 & 5.44 \\
  & Gemini 2.5 Flash$^*$            & 53.8 & 48.2 & 36.0 & 33.0 & 7.04 & 5.84 \\
  & Gemini 2.5 Flash+Think$^{\dagger*}$ & 71.6 & 74.4 & 58.0 & 57.0 & 4.64 & 4.72 \\
\midrule
\multirow{4}{*}{Qwen}
  & QwQ-Plus   & 72.8 & 72.6 & 59.0 & 58.0 & 4.80 & 4.72 \\
  & Qwen-Plus  & 63.4 & 61.2 & 52.0 & 51.0 & 4.72 & 4.64 \\
  & Qwen-Max   & 30.6 & 30.4 & 26.0 & 25.0 & 2.24 & 2.24 \\
  & Qwen-Turbo & 34.4 & 30.8 & 32.0 & 28.0 & 0.88 & 1.04 \\
\midrule
\multirow{2}{*}{DeepSeek}
  & DeepSeek-R1   & 84.2 & \textbf{84.4} & 74.0 & \textbf{75.0} & 3.76 & 3.20 \\
  & DeepSeek-V3.2 & 48.4 & 47.6 & 36.0 & 38.0 & 5.52 & 4.00 \\
\bottomrule
\end{tabular}
}
\end{table*}

\subsection{Overall Accuracy and Stability}
Table~\ref{tab:overall-performance} reports RLPR, PSR, and AV for all 32 model/prompt configurations. The top performers under RLPR are o4-mini (86.2\% D / 85.6\% M), Gemini 3.1 Pro (86.2\% D / 85.8\% M), and DeepSeek-R1 (84.2\% D / 84.4\% M). More broadly, the table shows that high run-level performance and high retry-free coverage often co-occur, but not in a one-to-one way: models with similar RLPR can still differ materially in PSR and AV.

Figure~\ref{fig:acc-stability-scatter} plots RLPR against PSR for all 32 configurations. All points fall below the diagonal, confirming that RLPR consistently overstates production reliability. Across all configurations, RLPR and PSR are strongly correlated ($r=0.985$, $p<0.001$), which is directionally expected because both metrics are derived from the same underlying per-problem success structure. The informative result is therefore not the existence of a high correlation itself, but the persistent and non-uniform gap between run-level success and retry-free coverage.

The gap $\Delta = \mathrm{RLPR} - \mathrm{PSR}$ is non-uniform and reaches its maximum in the mid-accuracy tier. Gemini 2.5 Flash has the largest gap (17.8\% under \textsc{Detailed}), followed by Gemini 3.1 Flash-Lite (16.2\% D) and Gemini 2.5 Flash+Think (17.4\% M). At the low end, Qwen-Turbo has a gap of only 2.4\%---not because it is uniformly reliable, but because its outcomes are highly polarized: many problems are either always failed or always passed, leaving less room for intermediate trial-to-trial variability. This pattern is reflected in AV: Qwen-Turbo has the lowest AV ($0.0088$), while Gemini 2.5 Flash has the highest ($0.0704$). Together, PSR and AV help distinguish between systems that are consistently strong, consistently weak, and unstable in the middle.

While it is mathematically expected that RLPR $\geq$ PSR for any non-trivial Bernoulli process, and that the two metrics are highly correlated overall, the \emph{non-uniformity} of $\Delta$ across models is the substantive empirical finding. This non-uniformity is itself structurally expected: Bernoulli variance is maximized at $p \approx 0.5$, so mid-performing models---whose per-problem success probabilities cluster in the intermediate range---will exhibit the largest gaps by construction. The practical contribution is that this gap is observable, model-specific, and meaningfully large (up to 17.8 percentage points), meaning that RLPR rankings alone cannot reveal how much of a model's accuracy translates into retry-free reliability.

Knowing the theoretical direction does not eliminate the need to measure the magnitude. The empirically determined quantities---how large $\Delta$ actually is (up to 17.8 pp) and which specific models fall in the mid-accuracy danger zone---cannot be derived from first principles, since per-problem success probabilities are not known a priori. The contribution is therefore best understood as empirical calibration of a theoretically expected pattern at scale.

\begin{figure}[t]
\centering
\includegraphics[width=0.92\linewidth]{figures/paper_figure1_accuracy_vs_stability.png}
\caption{RLPR vs.\ PSR scatter plot across all 32 configurations (16 models $\times$ 2 prompt types). Points are colored by model family and shaped by prompt type. The diagonal marks $\mathrm{PSR}=\mathrm{RLPR}$; points below indicate configurations where mean success overstates stable problem coverage.}
\label{fig:acc-stability-scatter}
\end{figure}

\subsection{Problem-Level Stability Analysis}

Figure~\ref{fig:stability-heatmap} presents a cross-model stability overview. Each row corresponds to one of the 16 models, ordered by RLPR under \textsc{Detailed} prompting (highest at top). Each column corresponds to one of the 100 problems, grouped by difficulty tier (Easy, Medium, Hard) and sorted within each tier by decreasing mean pass rate across all models. Cell color encodes the empirical pass probability for that model--problem pair, averaged over $R=5$ runs: green indicates a problem solved consistently on every run, red indicates consistent failure, and yellow indicates high trial-to-trial variability.

\begin{figure*}[t]
\centering
\includegraphics[width=0.96\linewidth]{figures/paper_figure2_stability_heatmap.png}
\caption{Cross-model stability heatmap (\textsc{Detailed} prompt). Rows: 16 models sorted by RLPR descending. Columns: 100 problems grouped by difficulty tier with dashed separators, and ordered within each tier by decreasing mean pass rate. Cell color encodes empirical pass rate over $R=5$ runs (green $=1.0$, red $=0.0$, yellow $\approx 0.5$).}
\label{fig:stability-heatmap}
\end{figure*}

Figure~\ref{fig:success-distribution} shows the distribution of per-problem pass rates aggregated across all models. Qwen-Turbo shows the most polarized profile (AV $= 0.0088$), while Gemini 2.5 Flash has the highest AV ($0.0704$), indicating substantial trial-to-trial variability.

This view is useful because it exposes qualitatively different reliability profiles that similar aggregate pass rates can hide. A model with many $\hat{p}_i \approx 1$ and $\hat{p}_i \approx 0$ problems behaves differently from a model with many $\hat{p}_i \approx 0.4$ or $0.6$ problems, even when their average run-level pass rate is similar.

\begin{figure}[t]
\centering
\includegraphics[width=0.92\linewidth]{figures/paper_figure3_success_distribution.png}
\caption{Problem-level success distribution (\textsc{Detailed} prompt). The figure contains one panel per model (16 panels total), arranged in decreasing order of PSR. Each panel shows the fraction of that model's 100 problems falling in each empirical pass-probability bin $\hat{p}_i \in \{0/5, 1/5, \ldots, 5/5\}$, where $0/5$ means the model failed all 5 runs and $5/5$ means it passed all 5.}
\label{fig:success-distribution}
\end{figure}

We further note that retry-based recovery does not eliminate instability. PSR directly measures the fraction of tasks that do not require retries, while AV highlights how much of the remaining workload lies in a stochastic regime.

\subsection{Difficulty Gradient and Failure Modes}
\label{sec:difficulty}

Pooling outcomes by benchmark difficulty reveals a pronounced gradient: run-level pass rates are highest on Easy problems (90.6\%), lower on Medium (62.2\%), and substantially lower on Hard (34.0\%). The dominant non-accept verdict is \emph{Wrong Answer} ($\sim$58\%), followed by \emph{Time Limit Exceeded} ($\sim$27\%) and \emph{Runtime Error} ($\sim$11\%), with the remaining $\sim$4\% comprising compile errors and other verdicts. Most failures thus arise from incorrect algorithms rather than formatting noise, though timeouts and runtime failures increase on harder problems, indicating that instability is also tied to implementation robustness. A low PSR can therefore reflect two regimes---algorithmic inconsistency and implementation fragility---which imply different mitigation strategies.

\subsection{Reasoning vs.\ Standard Models}

Reasoning-specialized models generate an internal chain-of-thought before producing a final answer, which could plausibly improve output consistency by structuring the solution process. We test whether this architectural distinction manifests as a systematically smaller accuracy--stability gap after controlling for overall accuracy level.

\paragraph{Grouping.} We treat four models as reasoning-specialized (marked $\dagger$ in Table~\ref{tab:models}): DeepSeek-R1, QwQ-Plus, o4-mini, and Gemini 2.5 Flash+Think. The remaining 12 systems are classified as standard models. All conclusions from this analysis remain exploratory given the modest group sizes.

\paragraph{RLPR-controlled gap analysis.} Because $\Delta = \mathrm{RLPR} - \mathrm{PSR}$ is itself correlated with RLPR (higher-accuracy models tend to have larger absolute gaps, as shown in Section~\ref{sec:gap}), a raw comparison of mean gaps between groups would be confounded by accuracy level. To address this, we fit a linear model $\Delta \sim \mathrm{RLPR}$ across all 32 configurations ($R^2=0.40$, $p<0.001$) and examine the residuals: a negative residual indicates a smaller gap than expected at the model's accuracy level, i.e., greater stability relative to accuracy.

Figure~\ref{fig:reasoning-scatter} shows the result. Reasoning models show strongly heterogeneous behavior: DeepSeek-R1 falls well below the regression line (mean residual $\approx -4.4$), suggesting markedly better stability than expected for its accuracy level; o4-mini is near the line (residual $\approx -0.5$); while QwQ-Plus (residual $\approx +1.7$) and Gemini 2.5 Flash+Think (residual $\approx +3.0$) both fall above it, showing larger-than-expected gaps. Across all eight reasoning-model configurations, the mean residual is $-0.04$, compared with $+0.01$ for the 24 standard-model configurations---a difference of 0.05 percentage points. A permutation test (10{,}000 iterations, two-sided) yields $p=0.97$, providing no evidence that the residual difference is systematic.

\begin{figure}[t]
\centering
\includegraphics[width=0.90\linewidth]{figures/paper_figure5_reasoning_vs_standard.png}
\caption{Accuracy--stability gap ($\Delta$) vs.\ RLPR for all 32 configurations. Reasoning models (diamonds) and standard models (circles) are shown against an OLS fit. Points below the line have smaller gaps than expected at their accuracy level.}
\label{fig:reasoning-scatter}
\end{figure}

\paragraph{Variance comparison.} Average variance (AV) provides a complementary view of trial-to-trial variability. The four reasoning models have a mean AV of 4.30, compared with 3.81 for the 12 standard models---again showing no systematic advantage for reasoning-specialized architectures on this metric.

\paragraph{Interpretation.} The four reasoning models do not exhibit systematically smaller accuracy--stability gaps or lower AV once accuracy is accounted for. DeepSeek-R1 shows suggestive evidence of above-expected stability (residual $\approx -4.4$\,pp), warranting verification in a larger sample. However, with only four reasoning models, the analysis is severely underpowered: the permutation test ($p=0.97$) should be read as ``insufficient evidence from a small sample'' rather than ``the effect is negligible.'' The pattern is model-specific: chain-of-thought reasoning neither reliably improves nor worsens stability within this sample. One important caveat: the Gemini 2.5 Flash vs.\ Flash+Think comparison is confounded by output token budget (4{,}096 vs.\ 24{,}576 tokens), so the pair should not be interpreted as a clean ablation of thinking mode alone.

\subsection{Prompt Sensitivity}

Figure~\ref{fig:prompt-comparison} visualizes the prompt effect as a paired scatter plot: each of the 16 models appears as a single point with its \textsc{Minimal}-prompt PSR on the $x$-axis and \textsc{Detailed}-prompt PSR on the $y$-axis. Points above the diagonal indicate that detailed prompting improves stability.

\begin{figure}[t]
\centering
\includegraphics[width=0.8\linewidth]{figures/paper_figure4_prompt_comparison.png}
\caption{Prompt sensitivity: paired scatter of PSR under \textsc{Detailed} vs.\ \textsc{Minimal} prompting for all 16 models. Points above the diagonal indicate that detailed prompting improves stability. Points are colored by model family.}
\label{fig:prompt-comparison}
\end{figure}

Prompt effects are model-dependent: most models show small differences (typically within 1--3 percentage points of PSR), but Claude Sonnet 4.5 shows a reversal (higher PSR under \textsc{Minimal}), while Gemini 3 Flash shows the opposite. No single template uniformly dominates---consistent with prior findings on prompt sensitivity \cite{zhuo2024prosa}. Exact McNemar tests on the paired PSR indicators find no model reaching $p<0.05$; the largest PSR shift is 6 points (o4-mini, whose decoding regime is not user-controllable; see the Limitations section), and among the remaining 15 models the largest shift is 4 points ($p=0.219$). Across 16 models, 11 favor \textsc{Detailed} and 5 favor \textsc{Minimal} (mean absolute difference $\sim$2.3\,pp). This is a deliberately low-contrast comparison---both templates share identical structure, differing only in instruction density---so the null result may partly reflect insufficient contrast rather than a true absence of prompt effects. With $N=100$, a 4-point shift has less than 20\% power, meaning meaningful effects cannot be ruled out at the current sample size.

\subsection{Accuracy--Stability Gap}
\label{sec:gap}

Although RLPR and PSR are strongly correlated overall ($r=0.985$), the gap $\Delta = \mathrm{RLPR} - \mathrm{PSR}$ is practically significant and non-uniform across models. Across all 32 configurations, $\Delta$ ranges from 2.4\% (Qwen-Turbo, \textsc{Detailed}) to 17.8\% (Gemini 2.5 Flash, \textsc{Detailed}). The gap peaks in the middle of the performance spectrum, where models succeed frequently enough to generate trial-to-trial variability but not consistently enough to eliminate it. This non-linearity means that a single RLPR ranking cannot be used to infer the magnitude of the accuracy--stability gap: two models with similar RLPR can differ substantially in how much of their accuracy converts to retry-free reliability. PSR makes this distinction explicit.

RLPR and PSR can disagree on model rankings despite the high correlation. For example, GPT-4.1 leads GPT-4.1-mini on RLPR (59.6\% vs.\ 58.4\%) but trails on PSR (47.5\% vs.\ 50.0\%)---a reversal driven by GPT-4.1's larger accuracy--stability gap (12.1\,pp vs.\ 8.4\,pp). This confirms that PSR captures deployment-relevant distinctions invisible to RLPR rankings alone.

\section{Conclusion}

We presented a repeated-run evaluation framework for deterministic programming tasks and used it to study the relationship between correctness and stability. Across 16 models, 100 problems, and 5 repeated runs per configuration (16{,}000 total instances), run-level pass rate and retry-free coverage are strongly correlated, but run-level accuracy consistently exceeds perfect stability. The accuracy--stability gap is non-uniform and reaches up to 17.8 percentage points, peaking in the mid-accuracy tier. Importantly, PSR and RLPR can reverse model rankings even at high correlation ($r=0.985$): GPT-4.1 leads GPT-4.1-mini on RLPR but trails on PSR, demonstrating that the two capture deployment-relevant distinctions invisible to single-run reporting. We view these metrics as complementary: RLPR characterizes average invocation success, while PSR and AV reveal whether that success is concentrated in consistently solved problems or dispersed across unstable outcomes. Prompt effects and the reasoning-vs-standard distinction are both heterogeneous and exploratory at the current sample size, suggesting that neither is a reliable lever for improving stability without model-specific tuning.

Future work should extend this framework in two directions: (1) replicating the repeated-run analysis on at least one non-coding deterministic NLP task (e.g., text-to-SQL) to validate generalizability, and (2) controlling output token budgets across standard and reasoning variants to cleanly isolate the effect of chain-of-thought reasoning on stability. While our evaluation isolates single-turn correctness, practical systems involve broader orchestration and runtime dependencies whose reliability and security properties require separate analysis \cite{jiang2026agenticaicybersecurityattack}. Taken together, these findings suggest that repeated-run evaluation is a necessary complement to standard accuracy reporting for any deterministic text-conditioned generation task in reproducibility-sensitive workflows.

\section*{Limitations}

Several limitations qualify our results. First, the benchmark is relatively small ($N=100$) and restricted to recent LeetCode-style problems. This design helps reduce direct overlap with heavily circulated historical benchmarks, but it does not eliminate contamination risk and should not be interpreted as proof of benchmark novelty. In particular, recently published contest problems are frequently discussed online within days of their release, meaning that recency alone is not sufficient to guarantee that evaluated models have not been exposed to equivalent problem statements through post-publication web text. The benchmark size is also a practical tradeoff: with 16 models, 2 prompt conditions, and 5 repeated runs per problem, even a 100-problem setup yields 16{,}000 judged evaluation instances. This scale is large enough to expose broad correctness--stability patterns, but still limited for detecting small prompt effects, which is consistent with the non-significant prompt comparisons reported in Section~4.5.

Second, our repeated-run protocol is intentionally narrow: we vary prompt template and repeated sampling while holding the evaluation harness fixed. As a result, the study does not characterize the effects of broader prompt engineering, tool use, self-repair, retrieval, test-time selection, or multi-turn interaction.

Third, model access is mediated through provider APIs and platform infrastructure. Even though task grading is deterministic, we do not control all backend implementation details, model refreshes, or service-side changes that may affect outputs over time. Our conclusions should therefore be interpreted relative to documented model snapshots and access dates rather than as timeless statements about any provider family.

Fourth, our summary metrics compress distinct failure modes into a binary pass/fail outcome. This is appropriate for end-task correctness, but it does not by itself explain whether instability arises from reasoning errors, formatting mistakes, runtime failures, or judge-specific edge cases.

Fifth, a specific methodological caveat applies to o4-mini. The o-series API does not expose temperature or top-$p$ controls; the model's internal decoding regime is not documented and may differ substantially from the $T=0.3$ setting used for all other models. If o4-mini's effective temperature is near zero---producing near-deterministic outputs across runs---then its measured PSR and AV reflect a qualitatively different quantity from the other systems: we would be observing the near-absence of stochasticity rather than meaningful stability under repeated stochastic sampling. This limits direct comparability of o4-mini's repeated-run metrics with those of other models, and any high PSR for o4-mini should be interpreted cautiously as a potential artifact of reduced output variance rather than as evidence of robust problem-solving reliability.

Finally, our conclusions are limited to Python code generation on deterministic coding tasks with executable tests. They should not be read as a complete account of reliability for other programming languages, for open-ended generation tasks, or for production systems that include additional orchestration and verification layers.

\section*{Ethics Statement}

This paper evaluates commercial and open-access LLM APIs on deterministic programming tasks. The study does not involve human subjects, personal data collection, or annotation labor. The main ethical considerations are instead tied to platform use and deployment interpretation. First, the evaluation relies on automated submission to an external online judge. Such use should be understood as research benchmarking under the platform's operational constraints, and any public release should avoid exposing credentials, session artifacts, or procedures that would enable abusive automation. Second, the evaluated models are accessed through third-party APIs whose outputs and availability may be governed by provider-specific terms of service and may change over time. Our reported results should therefore be interpreted as measurements of accessed systems under a documented evaluation protocol, not as immutable properties of any provider's model family.

More broadly, stronger performance on coding benchmarks can increase both beneficial and harmful capabilities. Better repeated-run reliability can support legitimate software engineering workflows, but it may also lower the barrier to generating code in settings where security review is weak. For this reason, we frame our contribution as an evaluation methodology rather than as an endorsement of fully autonomous code deployment.

\bibliography{references}

\appendix

\section{Prompt Templates}

\subsection{Detailed Prompt}

\begin{lstlisting}
<role>
You are an expert Python algorithm engineer solving deterministic
LeetCode-style problems.
</role>

<task>
Produce a correct Python solution that fits the provided template and
is optimized for the problem constraints.
</task>

<instructions>
1. Read the full problem statement, constraints, and template before
   writing code.
2. Produce the best practical algorithm for the stated constraints,
   aiming for the optimal time complexity when possible.
3. Do not return a brute-force, quadratic, exponential, or placeholder
   solution when the constraints require a more efficient algorithm.
4. Use the required method signature from the template exactly.
5. Ensure the algorithm is complete and submission-ready, not a sketch
   or partial attempt.
6. Handle edge cases implied by the statement and constraints.
7. Return one complete Python solution that can be submitted directly.
</instructions>

<format_requirements>
Return plain Python source code only.
Start directly with the code.
Do not include explanations, markdown fences, or surrounding commentary.
The code must be valid, runnable Python 3 with no syntax errors.
The code must fit the provided code template without requiring manual edits.
The final code should be written for Accepted-style performance, not just
sample-case correctness.
Only include comments when they clarify non-obvious logic.
</format_requirements>

<problem>
{problem_description}
</problem>

<code_template>
{code_template}
</code_template>
\end{lstlisting}

\subsection{Minimal Prompt}

\begin{lstlisting}
<role>
You are a Python coding assistant solving a deterministic
LeetCode-style problem.
</role>

<task>
Write one correct and efficient Python solution that fits the provided
template.
</task>

<instructions>
1. Infer an algorithm that matches the constraints before writing code.
2. Prefer the optimal or near-optimal solution for the input limits.
3. Do not output brute-force code unless the constraints clearly make
   it acceptable.
4. Use the exact method signature from the template.
5. Return only final submission code.
</instructions>

<format_requirements>
Return plain Python source code only.
Start directly with the code.
Do not include explanations, markdown fences, or extra text.
The code must be valid, runnable Python 3 with no syntax errors.
The code must fit the provided code template without requiring manual edits.
The solution must be efficient enough to avoid obvious Time Limit
Exceeded outcomes under the given constraints.
</format_requirements>

<problem>
{problem_description}
</problem>

<code_template>
{code_template}
</code_template>
\end{lstlisting}

\section{Model Snapshots}
\label{sec:model-snapshots}

\begin{table*}[t]
\centering
\small
\caption{Model snapshots used in the experiments. The API identifier column records the concrete model string used by the evaluation pipeline. Access dates are derived from the archived run dates for the corresponding completed paper runs.}
\label{tab:model-snapshots}
\resizebox{\textwidth}{!}{
\begin{tabular}{llll}
\toprule
Display Name & Provider & API Identifier & Access Date \\
\midrule
GPT-4.1 & OpenAI & gpt-4.1 & 2026-04-01 \\
GPT-4.1-mini & OpenAI & gpt-4.1-mini & 2026-03-31 \\
o4-mini & OpenAI & o4-mini & 2026-04-06 \\
Claude Sonnet 4.5 & Anthropic & claude-sonnet-4-5-20250929 & 2026-04-03 \\
Claude Haiku 4.5 & Anthropic & claude-haiku-4-5-20251001 & 2026-04-03 \\
Gemini 3.1 Pro (preview) & Google & gemini-3.1-pro-preview & 2026-03-28 \\
Gemini 3 Flash (preview) & Google & gemini-3-flash-preview & 2026-03-29 \\
Gemini 3.1 Flash-Lite (preview) & Google & gemini-3.1-flash-lite-preview & 2026-03-30 \\
Gemini 2.5 Flash & Google & gemini-2.5-flash & 2026-04-05 \\
Gemini 2.5 Flash+Think & Google & gemini-2.5-flash (thinkingBudget=8192) & 2026-04-05 \\
QwQ-Plus & Qwen & qwq-plus & 2026-03-23 \\
Qwen-Plus & Qwen & qwen-plus & 2026-03-29 \\
Qwen-Max & Qwen & qwen-max & 2026-03-23 \\
Qwen-Turbo & Qwen & qwen-turbo & 2026-03-23 \\
DeepSeek-R1 & DeepSeek & deepseek-r1 & 2026-03-24 \\
DeepSeek-V3.2 & DeepSeek & deepseek-v3.2 & 2026-03-25 \\
\bottomrule
\end{tabular}
}
\end{table*}

% -----------------------------------------------------------------------
% ACL / ARR Responsible NLP Research Checklist
% Required for ARR submissions (EMNLP 2025 / ACL 2025 cycle).
% Items are answered for:
%   "Accuracy, Stability, and Repeated-Run Reliability of Large Language
%    Models on Deterministic Programming Tasks"
% -----------------------------------------------------------------------
\section*{Checklist}

The checklist follows the guidelines from the ACL Responsible NLP Research
framework. For each question, the answer should be \textbf{Yes}, \textbf{No},
or \textbf{N/A}. The answer \textbf{Yes} means that the paper
provides the relevant information. If the answer is \textbf{No}, an
explanation must be given. If the answer is \textbf{N/A}, an explanation
may be given.

\begin{enumerate}

%  ------------------------------------------------------------------
\item \textbf{Did the paper include a section about limitations?}
  \textbf{Yes.} The paper contains a dedicated ``Limitations'' section
  that discusses benchmark size, contamination risk, API
  non-determinism, metric compression, o4-mini's undocumented decoding
  regime, prompt contrast limitations, and scope restrictions.

%  ------------------------------------------------------------------
\item \textbf{Did the paper include a section about ethical
  considerations?}
  \textbf{Yes.} The paper contains a dedicated ``Ethics Statement''
  section addressing platform-use norms, API terms of service,
  dual-use concerns around automated code generation, and the absence
  of human subjects or personal data.

%  ------------------------------------------------------------------
\item \textbf{Did you describe the limitations of your work?}
  \textbf{Yes.} Six distinct limitations are enumerated in the
  Limitations section: benchmark scope ($N=100$, one language),
  contamination risk, narrow experimental protocol, API opacity,
  metric compression (binary pass/fail), and o4-mini's uncontrolled
  decoding.

%  ------------------------------------------------------------------
\item \textbf{Did you discuss any potential negative societal impacts
  of your work?}
  \textbf{Yes.} The Ethics Statement acknowledges that improved
  repeated-run reliability could lower barriers to generating code in
  settings lacking adequate security review, and frames the contribution
  as an evaluation methodology rather than an endorsement of autonomous
  deployment.

%  ------------------------------------------------------------------
\item \textbf{Have you read the ethics review guidelines and ensured
  that your paper conforms to them?}
  \textbf{Yes.}

\end{enumerate}

\noindent\textbf{Artifacts}

\begin{enumerate}
\setcounter{enumi}{5}

%  ------------------------------------------------------------------
\item \textbf{Did you describe the study design, the creation of the
  evaluation benchmark, and the data collection process?}
  \textbf{Yes.} Section~3 (``Experimental Setup'') describes problem
  selection by recency rank, difficulty distribution, grading criteria,
  model families, decoding parameters, prompt templates, normalization,
  and execution pipeline.

%  ------------------------------------------------------------------
\item \textbf{Did you use existing data or create new data?}
  \textbf{Yes} (existing data). The evaluation benchmark consists of
  the 100 most recently published LeetCode problems at collection time.
  No new dataset was created; the problems and test cases are sourced
  from the LeetCode platform.

%  ------------------------------------------------------------------
\item \textbf{Did you discuss whether and how consent was obtained
  from people whose data you're using, if it applies?}
  \textbf{N/A.} The problems are algorithmically specified tasks
  containing no personal data. No personal identifiers or user-generated
  content beyond algorithmic problem statements are used.

%  ------------------------------------------------------------------
\item \textbf{Did you discuss whether the data you are using/curating
  contains personally identifiable information (PII)?}
  \textbf{Yes.} Section~3 explicitly states that the benchmark
  ``consist[s] of algorithmic specifications only and contain[s] no
  personally identifying information.''

%  ------------------------------------------------------------------
\item \textbf{Did you provide documentation of the artifacts, e.g.,
  coverage of domains and languages, demographic groups represented,
  etc.?}
  \textbf{Yes.} Section~3 documents benchmark composition: 100 problems,
  three difficulty tiers (20 Easy / 50 Medium / 30 Hard), algorithmic
  topics, Python-only evaluation, and selection methodology. Model
  snapshots (provider, API identifier, access date) are listed in
  Appendix~B.

%  ------------------------------------------------------------------
\item \textbf{Did you report the basic statistics of the data you used?}
  \textbf{Yes.} Benchmark statistics (size, difficulty split, topic
  coverage) are reported in Section~3.1. Per-model result distributions
  appear in Table~1, Figure~2, and Figure~3.

\end{enumerate}

\noindent\textbf{Computational experiments}

\begin{enumerate}
\setcounter{enumi}{11}

%  ------------------------------------------------------------------
\item \textbf{If you ran experiments, did you report the number of
  parameters for each model?}
  \textbf{No.} Parameter counts for the evaluated models are not
  publicly disclosed by all provider families (OpenAI, Anthropic, Google,
  Qwen, DeepSeek). The paper reports the API identifiers, model tier
  designations, and access dates in Table~1 and Appendix~B as the
  available model descriptors.

%  ------------------------------------------------------------------
\item \textbf{If you ran experiments, did you report the hyperparameters
  used and the range of hyperparameters you searched over?}
  \textbf{Yes.} Section~3.4 specifies temperature $T=0.3$, top-$p=0.9$,
  $R=5$ runs per problem, and the default 4{,}096-token output budget.
  Reasoning model budget exceptions are documented. The study does not
  perform hyperparameter search; values are fixed by experimental design.

%  ------------------------------------------------------------------
\item \textbf{Did you use the same evaluation procedure for all
  experiments?}
  \textbf{Yes.} All 32 model/prompt configurations are evaluated on
  the identical 100-problem set using the same deterministic judge,
  normalization pipeline, and metric definitions (RLPR, PSR, AV).

%  ------------------------------------------------------------------
\item \textbf{Did you report the total amount of compute used to run
  your experiments?}
  \textbf{No.} API-based inference costs depend on provider pricing and
  token consumption, which vary per model and are not uniformly
  available. The paper reports the number of evaluation instances
  (16{,}000 total) and the per-model output token budgets as proxies
  for compute scale. Exact compute cost in GPU-hours or dollar amounts
  is not reported because all inference is accessed through commercial
  APIs without access to underlying hardware metrics.

%  ------------------------------------------------------------------
\item \textbf{Did you use human annotators (e.g., crowdworkers) or
  research participants in any way?}
  \textbf{No.} All evaluation is automated. The judge is a deterministic
  code execution harness; no human annotation is used.

%  ------------------------------------------------------------------
\item \textbf{Did you report the instructions given to annotators and
  their payment/compensation?}
  \textbf{N/A.} No human annotators were used.

%  ------------------------------------------------------------------
\item \textbf{Did you discuss potential biases in the human annotation
  process and/or possible mitigation strategies?}
  \textbf{N/A.} No human annotation was used.

\end{enumerate}

\noindent\textbf{Model development (if applicable)}

\begin{enumerate}
\setcounter{enumi}{18}

%  ------------------------------------------------------------------
\item \textbf{Did you report model validation results?}
  \textbf{N/A.} The paper does not train or fine-tune models; it
  evaluates pre-existing commercial models through provider APIs.

%  ------------------------------------------------------------------
\item \textbf{Did you report the results of any ablation study?}
  \textbf{Yes.} A controlled prompt ablation (\textsc{Detailed} vs.\
  \textsc{Minimal}) is reported in Section~4.5 with paired McNemar
  tests. A reasoning-vs.-standard architecture comparison is reported
  in Section~4.4. Both analyses are explicitly scoped as exploratory
  given sample-size constraints.

%  ------------------------------------------------------------------
\item \textbf{Did you use the same evaluation procedure for all model
  variants?}
  \textbf{Yes.} Identical metrics, dataset, judge, normalization, and
  statistical inference procedures are applied to all 32
  model/prompt configurations (with the documented exception that
  o4-mini does not expose temperature/top-$p$ controls).

\end{enumerate}

\noindent\textbf{Computational social science / Information extraction}

\begin{enumerate}
\setcounter{enumi}{21}

%  ------------------------------------------------------------------
\item \textbf{If your work is on applications in the social sciences or
  other fields, did you discuss domain expertise?}
  \textbf{N/A.} The paper is an empirical NLP/ML evaluation study, not
  a social science or information-extraction application.

%  ------------------------------------------------------------------
\item \textbf{Did you describe potential applications and use cases of
  your work in a responsible manner?}
  \textbf{Yes.} The Conclusion frames the metrics as evaluation tools
  for reproducibility-sensitive software engineering workflows. The
  Ethics Statement explicitly cautions against using improved coding
  performance as justification for fully autonomous code deployment
  without security review.

\end{enumerate}

\end{document}
